% ch18-systems.tex -- An interpreter and a compiler in miniature
\chapter{An interpreter and a compiler in miniature}
\label{ch:systems18}
\index{interpreter}\index{compiler}\index{virtual machine}

The book opened with a program's journey from text to silicon and
closes by building the two endpoints small enough to hold in one
hand: a 9-opcode stack VM (the interpreter) and a code emitter for
the same expression (the compiler). Both are real, both run, and
the interpreter ships with a real off-by-one kept where you can
see it.

\begin{wins}[After this chapter you can]
Write a small compiler yourself: take a line such as
\texttt{2+3*4}, turn it into a tree, turn the tree into code for a
stack machine \emph{and} into RISC-V instructions, and run both to
check you get the same answer as a real C compiler. You can also read
a trace of a virtual machine, say exactly what each step did to the
stack, and put a measured number on what interpreting costs.
\end{wins}

\begin{rulebox}[Before you start]
\textbf{Recap in plain words.} A \emph{program} is a list of
instructions. The CPU you built in ch.\,5--8 (and, in miniature, the
8-bit CPU of Book~5 ch.\,16) runs such a list by repeating three
steps: \emph{fetch} the next instruction, \emph{decode} what it
means, \emph{execute} it. Ch.\,3 showed the RV32I instructions and ch.\,4
showed how functions use a stack. This chapter notices something
simple: nothing says the fetch-decode-execute loop has to be made of
silicon. A few lines of C can do the same job, and then we call it a
\emph{virtual machine} (VM) or \emph{interpreter}. A \emph{compiler}
is the other choice: instead of doing what the program says, it
translates the program into instructions for a real CPU once, and
then gets out of the way. Think of an interpreter as a translator who
sits beside you and repeats every sentence aloud, and a compiler as a
translator who hands you a finished book in your language.

\textbf{C in ten lines.} \texttt{int st[256]} is an array of 256 integers; here we use
it as a stack with an integer \texttt{sp} (the stack pointer) that
says how many items are on it. \texttt{st[sp++] = v} pushes
\texttt{v}; \texttt{st[--sp]} pops. \texttt{switch (x) \{ case 1: \dots; break; \}}
jumps to the \texttt{case} whose number equals \texttt{x}.
An \texttt{enum} gives names to small numbers. \texttt{while (1)} loops until a
\texttt{return}. \texttt{p[ip++]} reads the element at index
\texttt{ip}, then adds one to \texttt{ip}. Compile and run with
\texttt{gcc -O2 file.c \&\& ./a.out}.
\end{rulebox}

\begin{principle}[Think like this]
When a problem feels big, ask ``what is the smallest machine that
still shows the whole idea?'' A 9-instruction VM is not a toy to be
outgrown; it contains the same loop, the same stack and the same
decisions as the Java or Python machines that run most of the
world's software. Learn the small one completely and the big ones stop
being magic.
\end{principle}

\section{The divide: doing what the program says, or translating it}

Take a recipe written in French and a cook who only reads English. There are
two ways to get dinner. You can hire an interpreter who stands in the
kitchen, reads one line at a time, says it aloud in English, waits for the
cook to do it, and then reads the next line. Or you can hire a translator
who rewrites the whole recipe once, hands over an English copy, and leaves.
The first costs a little at every step, every time you cook. The second costs
a lot once and nothing afterwards. Those are exactly an \emph{interpreter}
and a \emph{compiler}.

The same split exists for programs. An interpreter reads the program as
\emph{data}, one instruction at a time, and carries out each one with its own
code. A compiler reads the program once and writes out a new program in the
machine's own instructions, which the CPU then runs directly. Python's
standard implementation interprets; C is normally compiled; Java and
JavaScript engines do a mix (the JIT section below).
\begin{figure}[htbp]
\centering
\begin{tabular}{lll}
\toprule
stage & interpreter & compiler \\
\midrule
read source & every run & once \\
decode & per instruction & per AST node \\
execute & dispatch loop & native code \\
cost & 10--100$\times$ & $\sim$1$\times$ \\
\bottomrule
\end{tabular}
\caption{The interpreter/compiler divide: the interpreter pays
decode-per-step forever; the compiler pays it once at build time.
Our VM ran 28 dispatched steps for 14 abstract operations.}
\label{fig:loop18}
\end{figure}


\begin{principle}[Concept: a program is just data]
The idea that unlocks this whole chapter is that an instruction list is a
\emph{value} that another program can read, check and transform. An
interpreter reads it and acts. A compiler reads it and writes out
another one. Once you see programs as data, building either is an exercise in
careful bookkeeping, not magic.
\end{principle}

\section{The interpreter: a stack machine}

\subsection*{What is a stack machine? Start from $2+3$}

A \emph{stack} is a pile where you can only add or remove at the top, like
plates. Suppose a calculator works on a stack. To compute $2+3$ you give it
three commands: \texttt{PUSH 2} (put 2 on the pile), \texttt{PUSH 3}, then
\texttt{ADD}. \texttt{ADD} takes the top two plates, adds them and puts
the result back. No command ever names where its operands are: they are always
``the top two plates''. That is why a stack machine's instructions are tiny
and why it is so easy to build one. Writing the operator after its
operands ($2\;3\;+$) is called \emph{postfix} notation, and it is exactly what
stack code is.

For a longer formula such as $(6\times7)+8$ the same idea works with no
brackets at all: push 6, push 7, \texttt{MUL} (leaves 42), push 8,
\texttt{ADD} (leaves 50). The order of the commands already carries the
structure that brackets expressed.
\begin{figure}[htbp]
\centering
\begin{tabular}{ll}
\toprule
op & stack after \\
\midrule
PUSH 6 & [6] \\
PUSH 7 & [6,7] \\
MUL & [42] \\
PUSH 8 & [42,8] \\
ADD & [50] \\
OUT & [] (prints 50) \\
\bottomrule
\end{tabular}
\caption{The stack machine's trace on \texttt{(6*7)+8}: operands
push, operators consume-and-push. No registers, no allocation ---
the simplest machine that computes.}
\label{fig:stack18}
\end{figure}


\texttt{ch18\_vm.c} implements 9 \emph{opcodes} (operation codes, the
numbers that name the instructions): PUSH, ADD, SUB, MUL,
DIV, DUP (copy the top of the stack), OUT (print and remove the top),
JNZ (jump if the top is not zero) and HALT, in one \emph{dispatch loop}.
Dispatch means: read the next opcode, pick the code that handles it,
run it, repeat. The program
computes $(6{\times}7)+8$, prints 50, then counts down via a JNZ
loop. The VM dispatches 28 instructions to do it (verified by running it,
and again by counting by hand below). Of those 28, 13 do real work
(5 arithmetic, 5 \texttt{OUT}, 3 \texttt{JNZ}) and 15 are stack
bookkeeping (7 \texttt{PUSH}, 7 \texttt{DUP}, 1 \texttt{HALT}): about
$28/13\approx 2.2$ dispatches for each useful operation. That ratio is
the interpreter's tax made visible, and it is a count of instructions,
not of time (the time cost is measured later in the chapter).

\begin{figure}[htbp]
\centering
\texttt{50 \quad 3 \quad 2 \quad 1 \quad 0 \quad (28 insn)}
\caption{Real \texttt{ch18\_vm} output verbatim: the arithmetic
program prints 50, then the countdown loop prints 3,2,1 ---
\emph{and a stray 0}, the JNZ exit landing one iteration late.
A real off-by-one kept as the chapter's receipt.}
\label{fig:vmrun18}
\end{figure}


\begin{example}[Why does the countdown print a stray zero?]
The countdown prints \texttt{3 2 1 0}, but a person reading the
description (``count down from 3'') expects \texttt{3 2 1}. Where does
the 0 come from? Read the program: inside the loop, each pass does
\texttt{DUP OUT} (print the counter), then \texttt{PUSH 1 SUB} (subtract
1), then \texttt{DUP JNZ} (jump back if the new value is not zero).
When the counter reaches 0 the jump is not taken, the loop is left,
and the next instruction is a final \texttt{OUT}, which prints the
value still sitting on the stack: 0. The VM did exactly what it was told;
the \emph{program} printed one value too many. This is the
loop-boundary mistake (an ``off-by-one'') that every hand-written
loop in every language can have. Keeping it in the chapter is
deliberate, because you will meet it again.

The fix is to delete that final \texttt{OUT}. The new bytecode
\texttt{ch18\_count\_fixed.bc} is \texttt{1 3 6 7 1 1 3 6 8 2 0} and
the loader \texttt{ch18\_vmrun} (below) prints, verbatim,
\begin{verbatim}
3
2
1
instructions executed: 20
\end{verbatim}
The count falls from 28 to 20 because the first part (the $6\times7+8$ arithmetic,
7 dispatches) and the final \texttt{OUT} (1) are gone.
\end{example}

\section{Watching the stack: a trace}
\begin{figure}[htbp]
\centering
\begin{tikzpicture}[x=1.7cm,y=0.55cm,font=\scriptsize]
\foreach \k/\lab in {0/{PUSH 6},1/{PUSH 7},2/{MUL},3/{PUSH 8},4/{ADD}}{
  \node[anchor=south] at (\k+0.5,2.4) {\lab};}
\foreach \k/\a in {0/6,1/6,2/42,3/42,4/50}{
  \draw[fill=accentsoft,draw=accent] (\k+0.1,0) rectangle (\k+0.9,1);
  \node at (\k+0.5,0.5) {\a};}
\foreach \k/\b in {1/7,3/8}{
  \draw[fill=goodbg,draw=good] (\k+0.1,1) rectangle (\k+0.9,2);
  \node at (\k+0.5,1.5) {\b};}
\node[anchor=east] at (0.05,0.5) {bottom};
\node[anchor=east] at (0.05,1.5) {next};
\node[anchor=north,text=slate] at (2.5,-0.2) {MUL pops 7 and 6, pushes 42; ADD pops 8 and 42, pushes 50};
\end{tikzpicture}
\caption{The operand stack of \texttt{ch18\_vmrun} after each step of
$6\times7+8$ (the real trace is in the worked example). Operands are
never named: the top of the stack is always ``the current value''.}
\label{fig:trace18}
\end{figure}


A stack machine is easiest to understand by watching it. The
loader \texttt{ch18\_vmrun.c} reads bytecode (a list of integers)
from standard input and, with \texttt{-t}, prints one line per
dispatched instruction showing the stack \emph{before} that step.
Unlike the first VM it also checks every push and pop, so a faulty
program stops with a message instead of silently overwriting memory.

\lstinputlisting[language=C,firstline=29,lastline=38,
caption={The dispatch loop of \texttt{ch18\_vmrun.c}: one \texttt{case}
per opcode, every pop and push checked.}]{code/ch18_vmrun.c}

\begin{example}[Following the stack step by step: a real trace of $6\times7+8$]
Feeding it the bytecode \texttt{1 6 1 7 4 1 8 2 0} (PUSH 6, PUSH 7,
MUL, PUSH 8, ADD, HALT) with \texttt{-t} gives, verbatim:
\begin{verbatim}
  1  ip=0   PUSH 6    stack before:
  2  ip=2   PUSH 7    stack before: 6
  3  ip=4   MUL       stack before: 6 7
  4  ip=5   PUSH 8    stack before: 42
  5  ip=7   ADD       stack before: 42 8
  6  ip=8   HALT      stack before: 50
instructions executed: 6
\end{verbatim}
Read it line by line. Step 3: the stack is \texttt{6 7}; \texttt{MUL}
pops the two top values, multiplies, pushes 42. Step 5: the stack is
\texttt{42 8}; \texttt{ADD} leaves 50. Notice that no instruction
says \emph{which} registers to use. That is why stack code is so
compact: the operands are always ``the top two things''.
\end{example}

\begin{example}[What if the program is wrong? The checked VM says so]
A VM that trusts its input is a security hole (the first VM, with
no checks, would let a long enough program write past
its 256-entry stack). The checked loader turns three classic
mistakes into messages. Verbatim:

{\scriptsize
\begin{verbatim}
PUSH 1, PUSH 0, DIV      ->  VM error at ip=5: divide by zero
(program: 6 6 6)          ->  VM error at ip=1: stack underflow
(300 times PUSH 1)       ->  VM error at ip=514: stack overflow
\end{verbatim}
}
The second program is \texttt{DUP DUP DUP}: the first \texttt{DUP}
tries to pop from an empty stack. The third pushes 300 values onto
a 256-entry stack and is stopped at the 257th push. The habit of mind: whenever a machine has a
fixed-size resource, ask ``what happens at the edge?'' and make the
answer a deliberate error, not an accident.
\end{example}

\section{The compiler: emit once}

The compiler's trade is to do the thinking \emph{before} the program runs. Our first
compiler, \texttt{ch18\_cc.c}, is as small as one can be: it
knows one expression, $(a+b)\times c$, and \emph{emits} (prints)
the instructions that compute it, as text.
\begin{figure}[htbp]
\centering
\begin{verbatim}
mov  w0, #6
mov  w1, #7
add  w0, w0, w1
mov  w1, #8
mul  w0, w0, w1
\end{verbatim}
\caption{\texttt{ch18\_cc} output: the same expression as 5 lines of
arm64 --- the compiler's advantage is that (a$+$b)$\times$c
becomes instructions \emph{once}; there is no dispatch in the
result.}
\label{fig:compile18}
\end{figure}


The output is five arm64 instructions with no dispatch loop and no stack
at run time. The constant operands are written straight into the
instructions, and the compiler even computes the answer itself
(\texttt{\# eval at compile time: 104}), a trick called \emph{constant
folding}. A plain interpreter, which looks at one instruction at a time as it
runs, does not do this unless someone adds a separate pre-pass,
and that pre-pass is already a small compiler.

\begin{principle}[Think like this]
Ask of any work: ``does this have to happen every time the program
runs, or can it happen once, earlier?'' Moving work from run time to build
time is the single idea behind most compiler optimisations, and later in this
chapter it is what the JIT does for the hot parts of a program.
\end{principle}


\section{A compiler you can read in one sitting}

So far the ``compiler'' only printed lines it already knew. A real
one has three jobs: \emph{parse} the text (find its structure), maybe
\emph{optimise} it, and \emph{generate} code. \texttt{ch18\_expr.py}
(about 120 lines of Python) does all three for integer expressions with
\texttt{+ - * /} and parentheses, and can aim at two targets: our stack
VM, or RISC-V assembly.

\paragraph{Why a tree?} The text \texttt{2+3*4} is a flat line of
characters, but its meaning has shape: the multiplication must happen
first, so \texttt{3*4} is a unit and the addition combines \texttt{2}
with that unit. A tree captures this: each node is an operator whose
children are its operands. The compiler's first job is to build the tree.

\begin{example}[First attempt: read left to right. It gets $2+3\times4$ wrong]
The simplest idea is to read left to right and apply each operator as
you meet it. Our \texttt{--naive} flag does exactly that, treating every operator as equal.
Its tree for \texttt{2+3*4} is $((2+3)\times 4)$ and its value is
$20$. The correct tree is $(2+(3\times4))$ with value $14$. Real output:
\begin{verbatim}
$ python3 ch18_expr.py --naive --target ast "2+3*4"
((2 + 3) * 4)
value: 20
$ python3 ch18_expr.py --target ast "2+3*4"
(2 + (3 * 4))
value: 14
\end{verbatim}
The cure is to give the grammar one level per precedence: an
\emph{expression} is terms joined by \texttt{+}/\texttt{-}; a
\emph{term} is factors joined by \texttt{*}/\texttt{/}; a \emph{factor}
is a number or a parenthesised expression. Because \texttt{expr} calls
\texttt{term}, which calls \texttt{factor}, the multiplication gets
grouped before the addition ever sees it.
\end{example}

\lstinputlisting[language=Python,firstline=42,lastline=56,
caption={The parser of \texttt{ch18\_expr.py}: a single helper
\texttt{binary} builds left-leaning trees for any level of operators.}]{code/ch18_expr.py}

\lstinputlisting[language=Python,firstline=81,lastline=92,
caption={Code generation. For the stack VM, walk the tree children-first
(``postorder'') and emit the operator last. For RISC-V, the same walk
hands out temporary registers.}]{code/ch18_expr.py}

\begin{example}[One expression, three kinds of output]
Compile $(6+7)\times 8$. Walking the tree children first gives the VM
bytecode, the register walk gives RISC-V, and with \texttt{--fold}
the compiler evaluates constants itself. Real outputs:

{\scriptsize
\begin{verbatim}
--target vmlist          --target rv32            --fold --target rv32
  0  PUSH 6                li   t0, 6               li   t0, 104
  2  PUSH 7                li   t1, 7               addi a0, t0, 0
  4  ADD                   add  t0, t0, t1
  5  PUSH 8                li   t1, 8
  7  MUL                   mul  t0, t0, t1
  8  OUT                   addi a0, t0, 0
  9  HALT
\end{verbatim}
}
Running the bytecode in \texttt{ch18\_vmrun} prints 104 after 7
dispatches. The RISC-V version needs no dispatch at all: 6
instructions, the last one just moving the answer into \texttt{a0},
the register that holds a function's result (ch.\,4). Assembled with
\texttt{clang --target=riscv32 -march=rv32im} and disassembled with
\texttt{llvm-objdump}, the multiply is the machine word
\texttt{b3 82 62 02}, which is \texttt{mul t0, t0, t1} (this needs the
M extension; RV32I alone has no multiply, as ch.\,3 noted). With
folding the whole program collapses to one \texttt{li}, which is
what \texttt{\# eval at compile time: 104} meant earlier in the chapter.
\end{example}

\begin{example}[How do we know the compiler is right? Compare it with gcc]
Correctness is checked, not assumed. \texttt{test\_ch18.sh} compiles
seven expressions, runs each on the VM, and compares with what gcc
computes for the same C expression (including C's rule that
integer division truncates). All seven match:
\texttt{104, 50, 5, 14, 7, 1, -3} for
\texttt{(6+7)*8, 6*7+8, 10-3-2, 2+3*4, 100/7/2, ((1+2)*(3+4))-5*(6-2), 7-10}.
Check one by hand: $100/7=14$ (truncated), $14/2=7$; and
$21-5\cdot4=1$. Programs using only \texttt{+} and \texttt{-}
also run on this book's own RV32I simulator from ch.\,3: the code for
$(9-4)+(7-5)$ leaves 7 in \texttt{a0} after 8 retired instructions.
\end{example}

\section{What does interpreting cost? Measure it}

Folklore says ``interpreters are 10 to 100 times slower''. Let us
see what our own VM does. \texttt{ch18\_vm\_bench.c} runs the same
countdown loop from 100 million, natively and as bytecode
(\texttt{PUSH N; L: PUSH 1, SUB, DUP, JNZ L; HALT}), four dispatches per iteration.

\lstinputlisting[language=C,firstline=13,lastline=21,
caption={The VM side of \texttt{ch18\_vm\_bench.c}: the same switch
dispatch as before, with no checks and no printing.}]{code/ch18_vm_bench.c}

\begin{example}[How much slower is interpreting? A fair baseline, then an unfair one]
Verbatim output (Linux x86\_64 container, \texttt{gcc -O2}):

{\scriptsize
\begin{verbatim}
iterations: 100000000
native, register counter:    72.3 ms  (0.72 ns/iter)
native, memory counter  :   168.5 ms  (1.69 ns/iter)
stack VM                :   716.6 ms  (7.17 ns/iter, 400000002
    dispatches, 1.79 ns/dispatch)
VM / native(register) = 9.9x    VM / native(memory) = 4.3x
\end{verbatim}
}
Two native baselines appear because the first one we tried, a
counter declared \texttt{volatile} so that the compiler could not delete the loop,
forces a store and reload on every pass and is itself slow
(1.69\,ns). Comparing against that flatters the VM ($4.3\times$). The honest baseline keeps the counter in a register,
which gives $9.9\times$. Over several runs the register-baseline ratio
ranged from 8 to 16$\times$ on this shared container, so quote it as ``about
10$\times$'', not as a precise figure. Each dispatch costs about 1.8\,ns, roughly 5
cycles: that is the price of loading the opcode, jumping to its
\texttt{case}, and touching the stack in memory. Folklore's 10--100$\times$ is therefore a
fair description of a naive interpreter on a loop with almost no real work in
it; real programs do more work per dispatch, which is why the figure
shrinks for them.
\end{example}

\section{The JIT compromise}

Compiling costs time up front, and most code in a typical program
runs only once or twice. A \emph{JIT} (just-in-time) compiler splits the
difference. It starts by interpreting everything, which costs nothing to
start. It counts how often each piece of code runs. When a piece
becomes \emph{hot} (it runs thousands of times), the JIT compiles just that
piece into machine code while the program is running. The cost of compiling is
then paid once and recovered many times over, like the translator hired only for
the chapters you re-read daily.

An optimising JIT goes one step further and \emph{speculates}: it
notices that a variable has always been an integer so far, compiles a
version that assumes it always will be, and puts a cheap check (a \emph{guard})
in front. If the guard ever fails, the engine abandons the fast code and
falls back to the interpreter, which is called \emph{deoptimising}. That is
why the figure has three tiers.
\begin{figure}[htbp]
\centering
\begin{tabular}{lll}
\toprule
tier & when it kicks in & cost model \\
\midrule
interpreter & always & decode-per-step \\
baseline JIT & warm code & compile once, simple code \\
optimising JIT & hot code & speculates + deopts \\
\bottomrule
\end{tabular}
\caption{The JIT answer to the divide: start interpreting, compile
the hot paths at runtime. Python's CPython lives at row one;
JVMs and JS engines climb all three.}
\label{fig:jit18}
\end{figure}


\section{The loop closes}

You can now trace one line of source text all the way down. Text is parsed into
a tree; the tree is translated to instructions (stack bytecode or RISC-V);
those instructions are fetched, decoded and executed, by a C loop
in the case of the VM or by the pipelined core of ch.\,5 to 8 in the case of
RISC-V; the core reads operands through the memory hierarchy of ch.\,9 and 10;
and the measurements of ch.\,15 tell you how fast it all was. Each layer you
met in this book is a place where one such translation happens.
\begin{figure}[htbp]
\centering
\begin{tabular}{ll}
\toprule
this book built & the layer it closed \\
\midrule
ch1 compile pipeline & source $\to$ binary \\
ch3--4 ISA/ABI & the contract \\
ch5--8 the core & the silicon \\
ch9--11 memory/OS & the environment \\
ch18 this machine & back to source \\
\bottomrule
\end{tabular}
\caption{The loop closes: source text in, computation out, through
every layer the book built. The 9-opcode VM is the same machine as
ch8's SoC --- one abstracts software, the other gates.}
\label{fig:full18}
\end{figure}


The VM is ch.\,8's SoC in software: fetch-decode-execute with a
\texttt{switch} for the decoder and a stack for the register file.
The compiler is ch.\,1's pipeline in miniature: syntax tree to instructions.
Everything between them --- the ISA contract, the memory
hierarchy, the OS services, the measurement discipline --- is what
this book spent 18 chapters building.

\section{Common mistakes}
\begin{itemize}
\item \textbf{``Interpreted = slow'' without a number.} Ours
      measured about 10$\times$ native on a bare countdown loop
      (4$\times$ against a deliberately slow native baseline) and 2.2
      dispatches per useful operation --- real, and dependent on how
      much work each dispatch does.
\item \textbf{The off-by-one at the loop boundary.} The printed
      trailing 0 is the \texttt{OUT} after the loop: test the exit
      condition and ask what is still on the stack when you leave.
\item \textbf{A parser that ignores precedence.} \texttt{2+3*4}
      came out as 20; one grammar level per precedence level fixes it.
\item \textbf{Thinking the compiler is magic.} Ours is about 120
      lines of Python. The hard parts (choosing registers when there are
      too few, serious optimisation) are real, but the \emph{idea} fits
      on a page.
\item \textbf{JIT as free speed.} A JIT pays to compile, needs warm-up
      time before the fast code exists, and loses its work whenever a
      guard fails and it must deoptimise. The interpreter pays a little all
      the time; the JIT pays more at the start and sometimes twice.
\item \textbf{Stack discipline bugs.} \texttt{OUT} removes what it
      prints, so \texttt{DUP} before \texttt{OUT} is what preserves the
      loop counter. One misplaced stack operation corrupts every later step.
\end{itemize}

\section{Summary}
\begin{tabularx}{\linewidth}{@{}l X@{}}
\toprule
\textbf{Piece} & \textbf{Verified version} \\
\midrule
VM       & 9 opcodes, real run: 50 + countdown 3..0 (the 0 is the program's, fixed version 3..1) \\
overhead & 28 dispatches, 13 of them useful work; $\approx$10$\times$ native on a bare loop (x86\_64) \\
compiler & same expr $\to$ 5 arm64 lines, no dispatch; \texttt{ch18\_expr.py}: parse, fold, emit VM or RV32 code \\
JIT      & interpret cold, compile hot, speculate \\
closure  & VM = ch8 SoC abstracted; book's loop closed \\
\bottomrule
\end{tabularx}

\section*{Further reading}
\begin{itemize}
\item Crafting Interpreters (Nystrom): the book of this chapter.
      \\ \texttt{craftinginterpreters.com}
\item The Python bytecode interpreter (ceval.c): \texttt{github.com/python/cpython}
\end{itemize}

\begin{rulebox}[Watch instead]
\begin{itemize}
\item \textbf{Crafting Interpreters overview --- Robert Nystrom
      talk} --- the author's walkthrough of this exact build
      (verified 2026-10-11).
      \url{https://youtu.be/8lC9FciP8Xg}
\end{itemize}
\end{rulebox}

\section*{Notes}
\begin{itemize}
\item The original programs run verbatim on this machine; the trailing 0 in
      the countdown comes from the program's final \texttt{OUT}, kept as
      a receipt. The new programs (\texttt{ch18\_vmrun.c},
      \texttt{ch18\_expr.py}, and \texttt{ch18\_vm\_bench.c}) are run and
      checked by \texttt{code/test\_ch18.sh}; timings are from a shared
      Linux x86\_64 container and vary by tens of percent between runs.
\item The earlier text spoke of ``14 abstract operations''; that figure
      cannot be derived from the program and was replaced by the counted
      13 useful and 15 bookkeeping dispatches.
\item The compiler's output is arm64 text emitted by the program ---
      the book ends where ch1 began, one level up.
\end{itemize}

\section*{Exercises}
\addcontentsline{toc}{section}{Exercises}
\stepcounter{section}

\begin{exercise}[easy]
Why can't the interpreter constant-fold like the compiler did?
\end{exercise}
\begin{solution}
The interpreter only ever looks at the instruction it is executing now;
it does not look ahead to notice that two pushed constants will be
multiplied. Folding needs a pass over the whole program before it runs, and a
pass that rewrites a program into a better one is compilation.
\end{solution}
\begin{exercise}[easy]
What does the 2$\times$ dispatch overhead buy the VM?
\end{exercise}
\begin{solution}
Portability and safety. The VM is ordinary C, so it runs on any machine that
has a C compiler, and it can check every operation (as \texttt{ch18\_vmrun}
does) instead of trusting the program. A compiler produces code for
one specific CPU and that code runs unchecked.
\end{solution}
\begin{exercise}[easy]
Which register does the stack machine never need to name?
\end{exercise}
\begin{solution}
All of them. Operands are implicit (``the top of the stack''), so
instructions carry no register fields. Register machines such as RV32I
must name each source and destination register.
\end{solution}
\begin{exercise}[easy]
Fix the countdown's trailing 0. One-line change?
\end{exercise}
\begin{solution}
Delete the final \texttt{OUT}: bytecode \texttt{1 3 6 7 1 1 3 6 8 2 0}
prints 3 2 1 in 20 dispatches (verified).
\end{solution}
\begin{exercise}[medium]
Add a CALL/RET pair to the VM. What stack discipline do the
caller and callee need that PUSH/ADD never did?
\end{exercise}
\begin{solution}
\texttt{CALL} must remember where to come back to, and \texttt{RET} must
find it. The caller leaves its arguments on the stack, \texttt{CALL} pushes
the return address (or pushes it on a second ``return stack''), and \texttt{RET}
pops it into \texttt{ip}. The callee must leave the stack as the caller
expects, for example only its result on top. That is the same
calling convention idea as ch.\,4's ABI, now in software.
\end{solution}
\begin{exercise}[medium]
The compiler baked (6+7)*8=104 at build time. What optimisation
pass is that, and when does it break (overflow, UB)?
\end{exercise}
\begin{solution}
It is \emph{constant folding}. It must give exactly the answer the program
would have produced at run time, so it has to respect the language's rules.
Our \texttt{ch18\_expr.py} stops with an error if a folded result does not
fit in 32 bits. In C, signed overflow is undefined behaviour, so a compiler is
allowed to fold however it likes; in languages with defined wrap-around it
must fold with wrap-around. Division by zero must also not be folded
silently (our compiler reports it).
\end{solution}
\begin{exercise}[medium]
Design the VM's dispatch differently: threaded code (jump table
per opcode). What does it remove and what does it cost?
\end{exercise}
\begin{solution}
A \texttt{switch} dispatch jumps from one shared place to the handler, so
all opcodes share one indirect jump that the predictor finds hard to guess.
Threaded code ends each handler with its own jump to the next handler, so
the predictor sees one jump per opcode with its own history. That removes
the shared bounds check and often mispredicts less. The costs are a less
portable source style (computed gotos, a GCC extension) and code that is harder
to read; the gain depends on the CPU and must be measured.
\end{solution}
\begin{exercise}[medium]
JIT speculation: your engine bets the loop counter stays int.
What must it emit to stay correct when it doesn't?
\end{exercise}
\begin{solution}
Before using the integer-specialised code it must check that the value
really is an integer (the guard). If not, it must jump to a \emph{deoptimisation}
path that rebuilds the interpreter's state from the compiled code's state
and carries on in the interpreter. Correctness always comes before speed.
\end{solution}
\begin{exercise}[hard]
Prove or refute: the stack VM's dispatch count per useful
operation grows with expression complexity. Count dispatches for an
expression with $n$ binary operators (compiled by \texttt{ch18\_expr.py})
and explain the scaling.
\end{exercise}
\begin{solution}
An expression with $n$ binary operators has $n+1$ numbers, so
$n+1$ \texttt{PUSH}es, $n$ operator instructions, one \texttt{OUT}
and one \texttt{HALT}: $2n+3$ dispatches (check: $(6+7)\times8$ has
$n=2$ and 7 dispatches). That is linear in $n$; per operator it tends to 2, so
the overhead does not grow.
\end{solution}
\begin{exercise}[hard]
Write the register-based VM version of this machine (reg[16]
instead of stack). Compare program byte length and dispatch count
on the same computation --- which machine wins and when?
\end{exercise}
\begin{solution}
Register code names its operands, so each instruction is bigger, but the
program needs fewer instructions because it has no \texttt{PUSH}/\texttt{DUP}
bookkeeping. For $(6+7)\times8$ the stack code has 7 dispatches and the
register code (see the RV32 output) 5 plus a final move into \texttt{a0}
(three constant loads and two arithmetic operations), so 6 instructions. Fewer dispatches generally mean a faster VM; the stack form is more compact.
The rule: stack code wins on size, register code on speed, which is why
modern VMs such as Lua's moved to registers.
\end{solution}
\begin{exercise}[easy]
What is the interpreter's \texttt{ip} in CPU terms?
\end{exercise}
\begin{solution}
The program counter. The CPU's PC holds the address of the next
instruction; the VM's \texttt{ip} holds the index of the next opcode.
Reading \texttt{p[ip++]} is fetch, the \texttt{switch} is decode, the case body is execute.
\end{solution}
\begin{exercise}[medium]
Why is the stack machine's code denser than a register machine's?
Give the byte-count argument on $(6{\times}7)+8$.
\end{exercise}
\begin{solution}
A stack \texttt{MUL} is just the opcode. A register
\texttt{mul t0,t0,t1} must also encode three register numbers (RV32I uses
32 bits per instruction). In our VM encoding $(6\times7)+8$ is
\texttt{PUSH 6, PUSH 7, MUL, PUSH 8, ADD}: if opcodes and constants
are one byte each, that is $2+2+1+2+1=8$ bytes, against $5\cdot4=20$ bytes for
five 4-byte RISC-V instructions.
\end{solution}
\begin{exercise}[medium]
The JIT tier's ``speculates and deopts'' --- what assumption does
a type-specialising JIT actually bet on?
\end{exercise}
\begin{solution}
That the types it has \emph{observed} so far will keep being the same,
for example ``this addition always receives integers''. The guard checks the bet
on every call, and a wrong bet costs a deoptimisation.
\end{solution}
\begin{exercise}[medium]
Our VM's stack lives in a C array. What breaks if a program
pushes past 256, and what's the honest fix?
\end{exercise}
\begin{solution}
In the unchecked \texttt{ch18\_vm.c} it silently writes past the array and
corrupts whatever memory is next to it, which is how real exploits work.
The honest fix is the one in \texttt{ch18\_vmrun}: check the stack pointer on
every push and pop and stop with an error, or grow the stack safely.
\end{solution}
\begin{exercise}[hard]
Design the minimal bytecode that a real JVM or WASM module would
need beyond our 9 ops to run ch12's scheduler sim. List the
semantic gaps (calls, locals, branches, memory).
\end{exercise}
\begin{solution}
Calls and returns with a frame for each call; local variables (load and store
by index); conditional and unconditional branches (not only JNZ), ideally structured
so they can be verified; and a block of addressable memory for arrays and structs.
These four, together with more integer types and comparison operations, are the
gaps between our toy and a real module such as WebAssembly.
\end{solution}
\begin{exercise}[easy]
Evaluate the bytecode \texttt{PUSH 2, PUSH 3, ADD, PUSH 4, MUL, OUT, HALT}
by hand, writing the stack after each step. What does it print, and
what infix expression is it?
\end{exercise}
\begin{solution}
Stacks: $[2]$, $[2,3]$, $[5]$, $[5,4]$, $[20]$; \texttt{OUT} prints 20.
It is $(2+3)\times4$, which is also what the naive parser computes for \texttt{2+3*4}.
\end{solution}
\begin{exercise}[medium]
Using the measured 7.17\,ns per iteration and 4 dispatches per
iteration, compute ns per dispatch, and convert to cycles at 2.8\,GHz.
\end{exercise}
\begin{solution}
$7.17/4=1.79$\,ns per dispatch; $1.79\times2.8\approx5.0$ cycles.
\end{solution}
\begin{exercise}[medium]
Compile \texttt{(1+2)*(3+4)} to RISC-V with the register scheme of
\texttt{gen\_rv}. How many temporary registers are live at the
busiest moment, and which?
\end{exercise}
\begin{solution}
The code is \texttt{li t0,1}, \texttt{li t1,2}, \texttt{add t0,t0,t1}, \texttt{li t1,3}, \texttt{li t2,4}, \texttt{add t1,t1,t2}, \texttt{mul t0,t0,t1}.
At the busiest point \texttt{t0}, \texttt{t1}, \texttt{t2} are live: three.
\end{solution}
